\documentclass[11pt,a4paper]{article}
\usepackage[utf8]{inputenc}
\usepackage[margin=2.4cm]{geometry}
\usepackage{parskip}
\usepackage[hidelinks]{hyperref}
\usepackage{enumitem}
\usepackage{xcolor}
\setlist[itemize]{leftmargin=1.4em, topsep=2pt}

\newcommand{\rc}[1]{\par\medskip\noindent\textbf{#1}\par\nopagebreak}
\newcommand{\cmt}[1]{\par\noindent\textit{#1}\par\smallskip}
\newcommand{\resp}{\par\noindent\textbf{Response.}\ }
\newcommand{\where}[1]{\par\smallskip\noindent\textcolor{blue!55!black}{\textbf{Where:} #1}\par}

\title{\vspace{-2.2cm}Response to Reviewers\\[4pt]
\large IJADS-312370 --- ``From Signal Fusion to Asset Allocation: A Decision-Theoretic
Framework for Regime-Conditional Portfolio Construction''}
\author{}
\date{}

\begin{document}
\maketitle
\vspace{-2.2cm}

\section*{Summary of the revision}

We thank both reviewers and the Editor for a careful and demanding review. Every
mandatory recommendation has been addressed, and all five optional suggestions have been
implemented as well. Page, table, figure and equation numbers below refer to the revised
manuscript.

One change is larger than a normal revision and we describe it plainly at the outset,
because the reviewers will notice it immediately.

Reviewer~2's third comment asked for robustness checks across rolling windows and for
sensitivity analysis. Implementing those checks exposed a defect in our own
implementation. The fused signal was being passed into the mean-variance objective in
the raw units of the underlying indicators; on-balance volume takes values of order
$10^{8}$ to $10^{9}$, whereas annualised variances are of order $10^{-2}$. The linear
term therefore dominated the quadratic risk term by roughly ten orders of magnitude, the
risk model had no influence on the solution, and the optimiser returned a corner
solution holding a single asset on almost every rebalance date. This is precisely the
behaviour Reviewer~2 identified in comment~4 as inconsistent with the described
risk-budgeting approach. The reviewer was right, and the cause was a defect rather than
a design choice.

A second defect surfaced from the same investigation. The regime-dependent volatility
targets were specified as absolute levels of 10 to 20~per cent annualised. On a
diversified portfolio those levels are never reached, so the constraint bound in one
rebalance out of 224 and the regime layer --- the central mechanism of the paper --- had
no effect on the portfolio at all.

We corrected both defects: expected returns are now scaled to a return metric
(Equation~(4), p.~10) and the regime risk budget is expressed as a fraction of the
portfolio's own volatility rather than as an absolute target (Equation~(7), p.~11).
Following Reviewer~2's optional comment~4, we also extended the evaluation from the
original five-year, ten-equity study to an eighteen-year, ten-instrument multi-asset
study covering four major stress episodes. The evaluation window runs to every session
available at the date the analysis was run, 26~August 2026, rather than stopping at the
original submission date. This is not a favourable window chosen after the fact: the
2008--2023 block reproduces exactly the figures the corrected pipeline produced before
the extension, and the 665 sessions from January~2024 onward postdate every
specification decision in the study. We therefore report them separately, in a new
Section~4.2, as an out-of-sample holdout. On that holdout the framework records a higher
Sharpe ratio and a smaller maximum drawdown than all three benchmarks, including the two
that beat it on Sharpe ratio before 2024.

The consequence is that all reported numbers have changed, and the conclusions are
narrower than in the submitted version. The framework's drawdown and stress-period
results survive and are now supported across 18 of 18 specifications; the claim of
significant alpha does not survive and has been removed; and the signal fusion layer is
now reported as a tested null, since ablation shows it contributes nothing. We judged
that reporting corrected results with weaker claims was the only defensible course, and
we would rather present this to the reviewers now than have it discovered later.

\bigskip
\hrule
\section*{Reviewer 1}

\subsection*{Changes which must be made before publication}

\rc{R1.1 Please fix your literature review. A typical academic paper should have around
45 academic references with about 10 relevant references published in ACADEMIC JOURNALS
in 2025 and 2026. Remove all self-citations and limit them to no more than 2--3
references. \dots\ no more than three citations should be attributed to the same journal
or author. In addition, please avoid citing references in clusters of more than three at
a time.}

\resp The reference list has been rebuilt from the ground up. It now contains
\textbf{49 references}, of which \textbf{27 were published in 2025 or 2026}, well beyond
the ten requested. Every entry was resolved against CrossRef, so authors, titles,
volume, issue, pages and DOI are taken from the publisher record rather than from
memory.

\begin{itemize}
\item \textbf{Journal concentration.} The submitted version drew 24 of its 33 references
from a single journal. No journal now supplies more than three: the maximum of three is
reached by The Journal of Finance, The Journal of Portfolio Management, IEEE Access and
this journal. Where a further article from this journal was worth citing, it replaced a
less relevant one rather than being added on top, so the limit of three is respected.
\item \textbf{Author concentration.} No author appears more than twice.
\item \textbf{Self-citations.} There are none, and there were none in the submitted
version.
\item \textbf{Citation clusters.} Every citation command in the manuscript was checked
programmatically; none groups more than three references.
\item \textbf{Scholarly sources.} All 49 entries are journal articles, apart from one
standard monograph (Grinold and Kahn, 2000) cited for the information-ratio scaling used
in Equation~(4). There are no conference proceedings, magazines or trade publications.
\item \textbf{IJADS references.} Following the Editor's request, three relevant articles
from this journal are now cited: Khodamoradi et~al. (2023) on robust
mean-conditional-value-at-risk portfolio optimisation, Zhang et~al. (2025) on the gap
between in-sample and out-of-sample model selection, and Mazzotta and Baidoo (2025) on
the informational value of alternative data.
\end{itemize}
\where{References, pp.~28--31; new citations throughout Section~2, pp.~4--7.}

\rc{R1.2 The paper has some serious acronym-related problems. \dots\ You need to expand
an acronym once in the abstract upon the first use and once in the body upon the first
use.}

\resp The abstract has been rewritten and now contains no unexpanded acronym. In the
body, every acronym is expanded at first use: compound annual growth rate (CAGR) and
maximum drawdown (Max DD) in Section~3.8; the Akaike and Bayesian information criteria
(AIC, BIC) in Section~3.4; information coefficient (IC) in Section~3.5; percentage
points (pp) in the caption to Table~6. The five technical indicators are defined in full
in Table~2 (p.~9), which is their first appearance anywhere in the manuscript. We also
reduced acronym use overall, writing ``maximum drawdown'' in running text and reserving
the abbreviation for table headings.
\where{Abstract, p.~1; Sections~3.4, 3.5 and 3.8, pp.~10--12; Table~2, p.~9.}

\rc{R1.3 Please carefully polish and clean up all figures and tables. All figures and
tables in a paper must be consistent. Please make sure all figures are high-resolution
(300 dpi).}

\resp All eight figures have been regenerated at exactly 300~dpi from the corrected
pipeline, in a single consistent visual style: one typeface, one colour palette, matched
line weights, matched axis treatment and matched figure widths. The framework diagram
(Figure~1) was redrawn from scratch, both for resolution and because the submitted
version depicted a news and text ingestion path that the implementation does not
contain.

Two further faults in the submitted version have been corrected. All seven figures
previously shared the same \verb|\label|, so every cross-reference resolved to the same
figure; each figure now carries a unique label and is referenced individually in the
text. All tables have been rebuilt to a single format.
\where{Figures~1--8, pp.~8, 14, 17, 20, 22; Tables~1--13 throughout.}

\bigskip
\hrule
\section*{Reviewer 2}

\subsection*{Changes which must be made before publication}

\rc{R2.1 The methodology should be explained more clearly so that other researchers can
understand and reproduce the work.}

\resp Section~3 has been rewritten with reproduction as its explicit objective, and now
runs from p.~7 to p.~12. It states the study universe instrument by instrument, the data
source and adjustment, the download window and the resulting evaluation period; gives
the exact formula and window for all nine input features in a single table; specifies
the mixture model and its refit protocol; states the fusion rule and the confidence term
as equations; gives the optimisation problem with every constraint; and sets out the
walk-forward protocol and the evaluation design. Nothing in the framework is now
described only in prose.
\where{Section~3, pp.~7--12.}

\rc{R2.2 More implementation details are needed, particularly regarding the data sources
used, sampling frequency, sentiment analysis model, feature extraction process, Gaussian
Mixture Model settings, and the optimization solver applied.}

\resp Each item is now specified explicitly.

\begin{itemize}
\item \textbf{Data source and sampling frequency.} Daily OHLCV series from Yahoo Finance
through the \texttt{yfinance} interface, adjusted for splits and dividends; download
window 1~January 2007 to 27~August 2026; evaluation period January~2008 to
August~2026, comprising 4{,}689 sessions, of which the final 665 are held out
(Section~3.1, p.~7).
\item \textbf{Sentiment model.} See R2.o3 below. In short: there is none, and the
manuscript now says so.
\item \textbf{Feature extraction.} Table~2 (p.~9) gives the closed-form definition and
rolling window of all nine features, followed by the cross-sectional standardisation in
Equation~(1) and the two composite definitions.
\item \textbf{Gaussian mixture settings.} Three components, full covariance matrices,
three random restarts, fixed seed, fitted on two standardised features, refitted at
every rebalance on an expanding window with a minimum of 252 observations
(Section~3.4, p.~10).
\item \textbf{Solver.} CVXPY with the Clarabel interior-point solver; covariance by
Ledoit--Wolf shrinkage on 252 trailing sessions (Section~3.6, p.~11).
\end{itemize}

We also documented one implementation detail that is easy to miss and that we ourselves
initially handled incorrectly: mixture component indices permute freely between refits,
so labels must be anchored --- here by sorting components on fitted volatility --- or the
regime series is silently scrambled (Section~3.4, p.~10).
\where{Section~3, pp.~7--12; Table~2, p.~9.}

\rc{R2.3 The experimental results also need stronger statistical support \dots\
statistical significance tests, confidence intervals, robustness checks using different
rolling windows, turnover statistics, and sensitivity analysis.}

\resp All five are now reported, and this comment is what led to the corrections
described in the summary above.

\begin{itemize}
\item \textbf{Significance tests.} Paired $t$-test and a Newey--West $t$-statistic with
ten lags on daily excess returns (Table~13, p.~23).
\item \textbf{Confidence intervals.} A stationary block bootstrap with 21-session blocks
and 2{,}000 replications gives interval estimates for the Sharpe ratio of the strategy,
of the benchmark, and of their difference (Table~13, p.~23).
\item \textbf{Robustness across windows.} A full sweep of three random seeds $\times$ two
covariance estimation windows $\times$ three rebalancing intervals. All eighteen cells
are reported, not a summary (Table~5, p.~16).
\item \textbf{Turnover statistics.} Annual one-way turnover, turnover per rebalance,
mean and minimum effective number of assets, mean largest weight, and mean and minimum
invested exposure (Table~11, p.~21).
\item \textbf{Sensitivity analysis.} Transaction costs varied from 0 to 25 basis points
(Table~11, p.~21), the regime risk budget varied over five specifications (Table~9,
p.~19), and component ablations at matched exposure (Table~8, p.~18).
\end{itemize}

We draw the reviewer's attention to what this analysis shows. The Sharpe improvement
over the benchmark is \emph{not} statistically significant: its bootstrap interval is
$[-0.102, 0.574]$ and includes zero. Section~4.11 (p.~23) says so directly, and the
abstract now states it too. What the evidence supports is consistency rather than
significance --- 18 of 18 specifications, 18 of 19 calendar years, 4 of 4 stress
episodes, and a holdout period the specification never saw --- and the manuscript is
careful not to conflate the two.
\where{Sections~4.3, 4.6, 4.7, 4.9 and 4.11; Tables~5, 8, 9, 11 and 13.}

\rc{R2.4 The portfolio allocation results require further clarification. In Figure~5,
the portfolio appears concentrated in only a few assets during certain periods, which
does not fully align with the described risk-budgeting approach. This behaviour should
be explained more clearly.}

\resp The reviewer identified a real defect, and we are grateful for it. The
concentration was not a property of the strategy but an artefact of unit mismatch in the
objective function, described in the summary above. In the submitted version the
portfolio held exactly one asset on nearly every date --- an effective number of assets
of 1.00, with a single holding persisting for 1{,}364 of 1{,}421 sessions --- because the
raw indicator scale overwhelmed the risk term by roughly ten orders of magnitude.

Equation~(4) (p.~10) now scales the fused score to a return metric before it enters the
objective, which resolves the mismatch. The corrected allocation is diversified and
stable: the effective number of assets averages 5.40 and never falls below 4.60, and
annual one-way turnover is 0.96 times portfolio value. Section~4.9 (pp.~20--22) explains
the failure mode explicitly rather than quietly replacing the figure, since the
diagnosis is itself a useful contribution.
\where{Equation~(4), p.~10; Section~4.9, pp.~20--22; Table~11, p.~21; Figures~5 and~6,
pp.~21--22.}

\rc{R2.5 The regime detection section also needs stronger justification. The choice of
three market regimes should be supported using model selection measures such as AIC,
BIC, or regime persistence statistics.}

\resp Table~12 (p.~23) now reports AIC and BIC for two through six components, together
with the persistence statistics for the three-state solution: the full transition
matrix, self-transition probabilities, mean run lengths, in-sample and out-of-sample
state shares, and the annualised volatility and momentum characterising each state.

We report an outcome that does not support our choice, and say so. Both AIC and BIC fall
monotonically across the range tested, so neither criterion selects three components;
this is the expected behaviour of a mixture fitted to persistent, heavy-tailed financial
features, which absorbs tail observations into additional components indefinitely. The
choice of three is therefore justified on persistence and interpretability instead, and
Table~12 supplies the grounds: the three states are economically distinct (annualised
volatility of 12.0, 20.3 and 49.4~per cent), they persist (self-transition probabilities
of 0.928 to 0.964, mean run lengths of 14 to 28 sessions), and the transition matrix is
close to tridiagonal, so the volatility ordering recovered by the model is economically
meaningful. Finer partitions produce states too short-lived for a monthly rebalancing
schedule to act on.
\where{Section~4.10, pp.~22--24; Table~12, p.~23; Figure~7, p.~24.}

\rc{R2.6 In addition, the language and presentation require careful revision. There are
grammatical inconsistencies, sudden changes in writing style, and some unfinished
formatting elements, including incomplete AI declaration text.}

\resp The manuscript has been rewritten end to end in a single voice, so the style
discontinuities are gone. Two specific faults the reviewer will have noticed are also
fixed: a stray sentence of reviewer-style commentary that had been left inside
Section~3.1.1 of the submitted version, and the misspelling of ``LLMs'' as ``LLPs'' in
Section~2.2.

The AI declaration placeholder has been replaced. We have also added a data and code
availability statement.
\where{Throughout; declarations on p.~29.}

\rc{R2.7 Finally, the equations should be explained in greater detail, with clear
definitions of all parameters and constraints used in the optimization model.}

\resp Section~3.6 (pp.~11--12) now presents the optimisation as two explicit steps and
defines every symbol. The objective terms and every constraint are given, along with the
values used and the reason for each: risk aversion $\lambda = 2$; position cap
$w_{\max} = 0.25$; turnover penalty $\kappa = c \cdot 252 / h$ with $c = 10$ basis points
and $h = 21$ sessions, annualised so that trading cost is on the same basis as return
and risk; the full-investment equality that makes the first step purely relative; and
the regime multiplier $k_r$ that sets total risk in the second step. Equations~(1)
through~(7) each carry an explanation of what the expression does and why it takes that
form, including why Ledoit--Wolf shrinkage is required and why the risk budget is
relative rather than absolute.
\where{Section~3.6, pp.~11--12; Equations~(1)--(7), pp.~9--12.}

\subsection*{Suggestions which would improve the quality of the article}

All five optional suggestions have been implemented.

\rc{R2.o1 The literature review can be strengthened by adding sharper critical
comparison between the proposed framework and recent portfolio allocation models.}

\resp Section~2.1 now compares the framework against mean--variance optimisation, naive
diversification, shrinkage-based optimisation, risk parity and hierarchical risk parity,
volatility management, regime-switching allocation and learning-based allocation.
Table~1 (p.~5) sets out each paradigm's basis of allocation, its principal limitation
and where this study sits, with our own limitation stated in the same terms as everyone
else's.
\where{Section~2, pp.~4--7; Table~1, p.~5.}

\rc{R2.o2 Figures would benefit from more technical interpretation, especially regime
transitions, signal confidence behaviour, allocation stability.}

\resp Every figure is now interpreted in the text rather than merely described.
Regime transitions are analysed through the transition matrix and run lengths
(Section~4.10), including the finding that direct calm-to-stress transitions have
estimated probability zero. Allocation stability is quantified --- effective number of
assets, largest weight, turnover --- and a new figure (Figure~6, p.~22) plots invested
exposure with the transitional and stress states shaded, so the de-risking mechanism is
directly visible. Signal behaviour is shown in Figure~8 (p.~24) and interpreted together
with the information coefficients in Table~10.
\where{Sections~4.1, 4.5, 4.9 and 4.10; Figures~2--8.}

\rc{R2.o3 The sentiment-analysis section may be improved by explicitly stating whether
domain-specific transformer models or lexicon-based methods were used.}

\resp Neither, and the revised manuscript states this plainly in a subsection written
for the purpose (Section~3.3, p.~9). No news article, social media post, corporate
announcement or any other text is read at any point in the study, and no lexicon,
classifier or language model is applied. What the submitted version described as a
sentiment pipeline is in fact a composite of four price and volume features, now
correctly named the market-state composite and defined in Table~2.

We went further than disclosure. Section~4.8 (pp.~19--20) measures the predictive
content of that composite directly and reports that it has none: its information
coefficient against next-session cross-sectional returns is $-0.026$
($t = -2.26$), and the regime-conditional weighting hypothesis --- that market-state
information becomes relatively more informative under stress --- is not supported: the
calm-to-stress contrast is insignificant ($t = -1.46$, $p = 0.145$) with point estimates
running the other way. Splitting at the holdout boundary also shows that the negative
coefficient is era-specific and disappears after 2023, so the layer cannot be rescued by
inverting its sign either. The corresponding portfolio-level ablation shows that removing
the layer entirely does not harm performance.
\where{Section~3.3, p.~9; Section~4.8, pp.~19--20; Table~10, p.~20; Table~8, p.~18.}

\rc{R2.o4 Future work may include cross-market validation, larger asset universes,
reinforcement learning-based adaptive parameter tuning.}

\resp Two of the three are now in the paper rather than deferred. The evaluation has
moved to a ten-instrument multi-asset universe spanning US, developed and emerging
equity, government and corporate fixed income, gold and real estate, over eighteen years
(Section~3.1, p.~7), and the original mega-capitalisation equity universe is retained as
a secondary study in Section~4.12 (p.~25). That secondary result is unfavourable and is
reported as such: the framework loses on risk-adjusted return in a concentrated
single-sector universe while retaining its drawdown advantage, which identifies the
condition under which the framework is useful.

Reinforcement-learning parameter tuning is retained as future work, now stated concretely
as learning the risk-budget multipliers, with the overfitting caveats that apply
(Section~5.5, p.~28).
\where{Section~3.1, p.~7; Section~4.12, p.~25; Table~14, p.~25; Section~5.5, p.~28.}

\rc{R2.o5 The manuscript would benefit from minor restructuring to reduce repetition in
Sections~4.2 and~4.3.}

\resp Those two sections have been merged. Section~4 is now organised so that each
subsection answers a distinct question --- headline performance, out-of-sample
holdout, robustness, stress
behaviour, annual consistency, ablation, budget form, signal content, allocation
behaviour, regime diagnostics, significance, secondary universe --- with no overlapping
material. Interpretation that previously appeared twice has been consolidated into
Section~5.
\where{Sections~4 and~5, pp.~13--28.}

\bigskip
\hrule
\section*{Changes not requested by the reviewers}

For completeness, three changes were made that no reviewer asked for.

\begin{itemize}
\item \textbf{Title.} ``Sentiment'' has been removed. Once Section~3.3 states that no
text or language model is used, the previous title asserted something the methodology
contradicts.
\item \textbf{Abstract.} The claim of ``significant alpha'' has been removed, because the
corrected results do not support it. The abstract now states the return sacrifice and
the absence of statistical significance alongside the positive findings.
\item \textbf{Contributions.} The list has been rewritten around what the evidence
supports: the two-step separation of composition from total risk, the matched-exposure
ablation design, the relative risk budget, and two documented failure modes.
\end{itemize}

We believe the revised manuscript is a substantially more careful piece of work than the
version submitted, and we thank the reviewers for the comments that forced it.

\end{document}
